\begin{abstract}
Este informe presenta la implementación desde cero, en Python/NumPy y
sin recurrir a funciones de filtrado adaptativo ya empaquetadas por
una librería externa, de dos algoritmos adaptativos clásicos: LMS
(\textit{Least Mean Squares}) en WORK\_ELEVEN y NLMS (\textit{Normalized
LMS}) en WORK\_TWELVE. Ambos se caracterizan sobre tres problemas
representativos del filtrado adaptativo: identificación ciega de un
sistema FIR desconocido a partir únicamente de pares entrada/salida,
cancelación adaptativa de ruido sobre un tono sintético de 60~Hz y
sobre una grabación de audio real, y el análisis cuantitativo de la
convergencia frente al tamaño de paso del algoritmo. Adicionalmente,
WORK\_TWELVE incorpora una comparación directa y cuantitativa entre
NLMS y una reimplementación local de LMS bajo una señal de entrada cuya
potencia cambia abruptamente a mitad del experimento, cuantificando en
aproximadamente 34 veces la ventaja de error de NLMS en el segmento de
menor potencia y evidenciando así la razón estructural —la
normalización del paso por la energía instantánea de la entrada— por
la cual NLMS resulta preferible a LMS en aplicaciones reales sobre
señales no estacionarias, como la voz.
\end{abstract}
